\documentclass[11pt,a4paper]{article}
\usepackage[utf8]{inputenc}
\usepackage[T1]{fontenc}
\usepackage{amsfonts}
\usepackage[french]{babel}
\frenchbsetup{StandardLists=true} % à inclure si on utilise \usepackage[francais]{babel}
\usepackage{enumitem}
\usepackage{amssymb}
\usepackage{amsmath}
\usepackage[left=2cm,right=2cm,top=2cm,bottom=2cm]{geometry}
\usepackage{multirow}
\usepackage[dvipsnames]{xcolor}

\usepackage[T1]{fontenc}
\usepackage{fouriernc}
\usepackage{graphicx}
\usepackage{colortbl}

\usepackage{setspace}
\singlespacing

\usepackage{pstricks}
\usepackage{xkeyval}
\usepackage{pst-labo}


\usepackage{multicol}
\usepackage{caption}
\usepackage{fancyhdr}
\pagestyle{fancy}
\renewcommand\headrulewidth{1pt}
\fancyhead[L]{Lycée Denis Diderot }
\fancyhead[R]{Classe de 1 Bac Pro }
\setcounter{page}{6}\fancyfoot[C]{page \thepage}
\fancyfoot[R]{Année 2024-2025}
\fancyfoot[L]{Alexis RUIZ}
\usepackage{multirow,tabularx,booktabs}
\newcommand{\cadreblanc}[1]{%
\medskip\framebox[\linewidth][c]{%
\rule{0mm}{#1mm}}\par
\medskip}

\usepackage{awesomebox}
\setlength{\aweboxvskip}{0mm}
\usepackage{pifont}
\usepackage{chemmacros}
\chemsetup{modules=all}
\setlength{\columnseprule}{.25mm}

\renewcommand{\arraystretch}{2}
\newcolumntype{C}[1]{>{\centering\arraybackslash }b{#1}}


\newcommand{\code}{\awesomebox[black]{2pt}{

\includegraphics[scale=0.15]{Chapitre 1 - Energie et puissance électrique/images/frame.png}
}{black}}



\newcommand{\defconclusion}{\awesomebox[black]{2pt}{\rotatebox{90}{\Large{\textbf{Conclusion}}}}{black}}


  \newenvironment{itemize*}%  % listes compactes
    {\begin{itemize}%
      \setlength{\itemsep}{0pt}%
      \setlength{\parskip}{0pt}}%
    {\end{itemize}}


\frenchbsetup{StandardLists=true}
\newlist{fleche}{itemize}{1}
\setlist[fleche]{label=\ding{220},font=\color{orange}}


\usepackage[tikz]{bclogo} 

\newcommand{\ctext}[1]
{\begin{bclogo}[couleur = white,logo= , arrondi = 0.1,barre=none]
{}
\vspace{-8mm}
~~
\textcolor{red}{#1}
~~
\end{bclogo}}

%%%%%%%%%%%%%%%%%%%%%%%%%%%%%%%%%%%%%%%%%%%%%%%%%%%%%%%%%%%



\begin{document}


\begin{center}
 \LARGE{COURS ET EXERCICES} \\[3mm]
\end{center}


\normalsize

\hrule\
\\[0,5cm]


\fcolorbox{black}{white}{\textbf{\ding{202} Pression et force pressante}}\\
\vspace{-3mm}

\begin{fleche}
\item \textbf{Pression }
\begin{itemize}[label=\textbullet, font=\LARGE \color{black}]
    \item \cadreblanc{10}
    \item \cadreblanc{30}
    \item \cadreblanc{20}
\end{itemize}
\vfill 
\item \textbf{Mesure de la pression}
\begin{itemize}[label=\textbullet, font=\LARGE \color{black}]
    \item \cadreblanc{10}
    \item \cadreblanc{10}
\end{itemize}


\vfill 

\item \textbf{Relation entre pression, surface et force pressante}

\noindent%
\begin{minipage}{.70\textwidth}%
\cadreblanc{10}
\end{minipage}%
\hfill
\begin{minipage}{.23\textwidth}%
\cadreblanc{10}
\end{minipage}%

\vfill 

\item \textbf{Pression absolue et pression relative}

\noindent%
\begin{minipage}{.7\textwidth}%
\cadreblanc{10}
\end{minipage}%
\hfill
\begin{minipage}{.23\textwidth}%
\fbox{\includegraphics[width=\textwidth]{Chapitre 3 - Pression et forces préssantes/images/cours1.png}}
\end{minipage}%

\newpage

\fcolorbox{black}{white}{\textbf{\ding{203} Loi de Boyle-Mariotte}}\\
\vspace{-3mm}


\begin{itemize}[label=\textbullet, font=\LARGE \color{black}]
    \item 
\noindent%
\begin{minipage}{.65\textwidth}%
\cadreblanc{10}
\end{minipage}%
\hfill
\begin{minipage}{.23\textwidth}%
\cadreblanc{10}
\end{minipage}%

    \item 
    \noindent%
\begin{minipage}{.90\textwidth}%
\cadreblanc{10}
\end{minipage}%
\end{itemize}
\noindent%
\begin{minipage}{.38\textwidth}%
\fbox{\includegraphics[width=\textwidth]{Chapitre 3 - Pression et forces préssantes/images/cours2.png}}
\end{minipage}%
\hfill
\begin{minipage}{.38\textwidth}%
\fbox{\includegraphics[width=\textwidth]{Chapitre 3 - Pression et forces préssantes/images/cours3.png}}
\end{minipage}%


\fbox{\textbf{Exercice 1 }} QCM : Parmi les trois propositions, choisissez la bonne réponse. 

\begin{center}



    \begin{tabular}{|C{0.5cm} p{5.5cm}|C{3cm}|C{3cm}|C{3cm}|}
         \hline
       \ding{202} &L'unité du système international de la pression est & $\Box $ ~~le bar& $\Box $ ~~l'atmosphère & $\Box $ ~~le pascal\\
         \hline
      \ding{203}&Quelle est la relation entre la pression et la force pressante  & $\Box $ ~~$p=F \times S$ & $\Box $ ~~$F=p \times S$  & $\Box $ ~~ $F=\dfrac{S}{P}$ \\
         \hline
         \ding{204} &Pour une même surface, si la force pressante augmente  que fait la pression &$\Box $ ~~ elle diminue& $\Box $ ~~ elle augmente & $\Box $ elle reste constante  \\
         \hline
         \ding{205} &Pour une même force pressante que peut on faire pour augmenter une pression & $\Box $ ~~ diminuer la surface de contact& $\Box $ augmenter la surface de contact& $\Box $ ~~ on ne peut pas la modifier\\
         \hline
         \ding{206}& La pression atmosphérique est environ égale à  &$\Box $ ~~ $1\times 10⁵$~Pa& $\Box $ $1\times 10^{-5} $~Pa& $\Box $ ~~ e$1\times 10⁵$~Bar \\
         \hline
         \ding{207}&Quel appareil ne permet pas de mesurer la pression & $\Box $ ~~ un baromètre & $\Box $ ~~  un manomètre & $\Box $ ~~  un anémomètre \\
         \hline
         %\ding{208} Un mouvement rectiligne uniformément ralenti peut être représenté par ...  & $\Box $ ~~  & $\Box $ ~~  &  $\Box $\\
         %\hline
         \ding{208} &Pour une même quantité de gaz, si la pression augmente que fait le volume&$\Box$ ~~ il diminue& $\Box$ il augmente &$\Box$ ~~ il reste inchangé\\
         \hline
         %\ding{209}&Dans un transformateur parfait, on a ...  &$\Box ~~ \dfrac{U_1}{U_2}=\dfrac{I_2}{I_1}$ &$\Box ~~ \dfrac{U_2}{U_1}=\dfrac{I_2}{I_1}$  & $\Box ~~ \dfrac{U_2}{I_1}=\dfrac{I_2}{U_1}$ \\ 
         %\hline 
         %\ding{210} &Le symbole électrique du transformateur est   &$\Box ~~ $ \includegraphics[scale=0.55]{Chapitre 2 - Transporter l'énergie sous forme électrique/images/QCM1.png}   &$\Box ~~ $ \includegraphics[scale=0.75]{Chapitre 2 - Transporter l'énergie sous forme électrique/images/QCM2.png}     &$\Box ~~ $ \includegraphics[scale=0.65]{Chapitre 2 - Transporter l'énergie sous forme électrique/images/QCM3.png}   \\
         %\hline

         
        
    \end{tabular}
\end{center}

\newpage

\fbox{\textbf{Exercice 2 - compléter le tableau suivant : }}

\begin{center}

    \renewcommand{\arraystretch}{3}


    \begin{tabular}{|p{4cm}|C{3cm}|C{3cm}|C{3cm}|}
         \hline
       & $P$ & $F$ & $S$ \\
         \hline 
               Formule à  utiliser & &  & \\
         \hline 
         Exemple 1 &  1,013 bar  & & 0,5 $m^2$\\
         \hline
         Exemple 2 &  & 50 N & 6,4 $cm^2$ \\
         \hline 
         Exemple 3 & $2 \times 10^5$ Pa & 50 N & \\
         \hline
         

    \end{tabular}\end{center}

    \vfill 


\fbox{\textbf{Exercice 3 }}

\vspace{-3mm}

\noindent%
\begin{minipage}{.36\textwidth}%
Un apnéiste inspire à la surface 6 L d'air à la pression atmosphérique. \'A l'aide de la loi de Boyle-Mariotte, calculer en L, le volume occuper par cet air à 20 m de profondeur, sachant que la pression augmente de 1 bar tous les 10 m.  

\end{minipage}%
\hfill
\begin{minipage}{.53\textwidth}%
\cadreblanc{40}
\end{minipage}%

\vfill 

\vfill 


\fbox{\textbf{Exercice 4 }}

Un ballon de baudruche peu gonflé et fermé est posé sous une cloche à vide. 

\begin{enumerate}
    \item Pourquoi ce ballon augmente-t-il de volume quand on chasse l'air sous la cloche à vide ? \\[2mm]
    \cadreblanc{15}
    \item Quelle doit être la pression sous la cloche pour que le ballon triple de volume ? \\[2mm]
    \cadreblanc{20}
\end{enumerate}


\vfill 

\end{document}